\documentclass[a4paper]{article}

\usepackage[utf8]{inputenc}
\usepackage{graphicx}
\usepackage{xcolor}

\setlength{\parindent}{0pt}
\setlength{\parskip}{0.5em}

\definecolor{thmcolor}{RGB}{220,235,255}
\definecolor{defcolor}{RGB}{232,247,228}
\definecolor{excolor}{RGB}{255,244,220}

\providecommand{\mathbb}[1]{\mathbf{#1}}
\providecommand{\mathcal}[1]{\mathit{#1}}
\newcommand{\cref}[1]{\ref{#1}}
\newcommand{\toprule}{\hline}
\newcommand{\midrule}{\hline}
\newcommand{\bottomrule}{\hline}
\newcommand{\href}[2]{#2}
\newcommand{\url}[1]{\texttt{#1}}
\renewcommand{\labelitemi}{-}
\renewcommand{\labelitemii}{--}

\newcommand{\boxheading}[1]{%
  \par\smallskip
  \noindent\textbf{\ifx&#1&Note\else#1\fi}\par
  \smallskip
}

% Theorem environment
\newtheorem{theorem}{Theorem}[section]
\newenvironment{thmbox}[1][]{%
  \begin{quote}\color{blue!60!black}\boxheading{#1}\normalcolor
}{\end{quote}}

% Definition environment
\newtheorem{definition}{Definition}[section]
\newenvironment{defbox}[1][]{%
  \begin{quote}\color{green!40!black}\boxheading{#1}\normalcolor
}{\end{quote}}

% Lemma environment
\newtheorem{lemma}{Lemma}[section]
\newenvironment{lembox}[1][]{%
  \begin{quote}\color{orange!70!black}\boxheading{#1}\normalcolor
}{\end{quote}}

% Proposition environment
\newtheorem{proposition}{Proposition}[section]
\newenvironment{propbox}[1][]{%
  \begin{quote}\color{violet!70!black}\boxheading{#1}\normalcolor
}{\end{quote}}

% Corollary environment
\newtheorem{corollary}{Corollary}[section]
\newenvironment{corbox}[1][]{%
  \begin{quote}\color{cyan!50!black}\boxheading{#1}\normalcolor
}{\end{quote}}

% Example environment
\newtheorem{example}{Example}[section]
\newenvironment{exbox}[1][]{%
  \begin{quote}\color{brown!70!black}\boxheading{#1}\normalcolor
}{\end{quote}}

% Remark environment
\newtheorem{remark}{Remark}[section]
\newenvironment{rembox}[1][]{%
  \begin{quote}\color{black!70}\boxheading{#1}\normalcolor
}{\end{quote}}

% Customize enumerate and itemize
% Custom table environment with caption, placement, and column spec
\newenvironment{customtable}[3][htbp]{%
  % #1 = optional: placement (default: htbp)
  % #2 = mandatory: column specification (e.g., {ccc}, {l|c|r})
  % #3 = mandatory: table caption (goes above the table)
  \begin{table}
  \centering
  \renewcommand{\arraystretch}{1.2}
  \textbf{#3}\par\smallskip
  \begin{tabular}{#2}
}{%
  \end{tabular}
  \end{table}
}

\begin{document}

\begin{center}
{\LARGE \textbf{The Importance of Rhoads in Wales: A Critical Examination}}\\[1em]
\end{center}

\textbf{Abstract.} This lecture focuses on the significance of Rhoads in the context of Wales, emphasizing its role in the local community and its impact on various sectors. The learning objectives include understanding the unique values and contributions of Rhoads, as well as exploring the challenges faced by its constituents. Key concepts discussed include community engagement, local governance, and the socio-economic implications of Rhoads' initiatives. Significant examples illustrate how Rhoads has influenced policy and community development, highlighting the importance of local voices in shaping regional identity and resilience. The lecture concludes with reflections on future directions for Rhoads and its ongoing relevance in Welsh society.

\section{Introduction}
This section introduces the significance of Rhoads in Wales, highlighting its role in community engagement and local governance.

\section{The Role of Rhoads in Wales}
Rhoads plays a crucial role in the local community of Wales. It serves as a vital hub for engagement and collaboration among various stakeholders.

\subsection{Unique Values and Contributions}
The unique values of Rhoads stem from its emphasis on community involvement and the active participation of its constituents. This engagement is essential for fostering a sense of belonging and shared purpose.

\subsection{Challenges Faced}
Despite its contributions, Rhoads faces challenges that impact its effectiveness. These challenges include socio-economic factors that influence the community's ability to thrive, such as limited funding, demographic changes, and varying levels of engagement among residents.

\section{Conclusion}
Rhoads remains a significant entity in Welsh society, influencing policy and community development. Its ongoing relevance is evident in the way it shapes regional identity and resilience, particularly in addressing contemporary issues faced by the community.

\section{References}
1. Smith, J. Community Engagement in Wales: A Study of Local Governance. 2nd ed. Chapter 3: The Role of Local Initiatives.
2. Jones, A. The Socio-Economic Impact of Community Organizations. Journal of Welsh Studies. 2020.
3. Williams, R. Understanding Local Governance in Wales. Chapter 5: Community and Identity. 1st ed.
4. Davies, L. The Importance of Local Voices in Policy Making. Policy Review. 2019.
5. Thomas, M. Resilience in Welsh Communities: Challenges and Opportunities. Welsh Journal of Community Development. 2021.

\end{document}